% !TeX root = ../main.tex
% Schemat blokowy architektury oprogramowania i przeplywu danych (Rys. 3.1)
\begin{figure}[H]
\centering
\resizebox{\textwidth}{!}{%
\begin{tikzpicture}[
    font=\sffamily\footnotesize,
    >=Stealth,
    node distance=5mm and 6mm,
    % Style warstw (ramki zewnetrzne)
    ctrlframe/.style={
        draw=black!55, fill=black!4, rounded corners=3pt,
        line width=0.8pt, inner sep=6pt
    },
    stageframe/.style={
        draw=blue!55!black, fill=blue!3, rounded corners=3pt,
        line width=0.9pt, inner sep=7pt
    },
    % Style blokow wewnetrznych
    ctrlbox/.style={
        draw=black!50, fill=white, rounded corners=2pt,
        line width=0.6pt, align=center, inner sep=4pt,
        minimum height=10.5mm
    },
    libbox/.style={
        draw=blue!50!black, fill=white, rounded corners=2pt,
        line width=0.65pt, align=center, inner sep=4.5pt,
        minimum height=10mm
    },
    tbbox/.style={
        draw=orange!75!black, fill=orange!10, rounded corners=2pt,
        line width=0.75pt, dashed, align=center, inner sep=4.5pt,
        minimum height=10mm
    },
    stagehdr/.style={
        font=\sffamily\footnotesize\bfseries, text=blue!55!black,
        anchor=west, inner sep=1pt
    },
    ctrlhdr/.style={
        font=\sffamily\footnotesize\bfseries, text=black!75,
        anchor=west, inner sep=1pt
    },
    dataflow/.style={
        ->, line width=0.95pt, draw=blue!55!black
    },
    internalflow/.style={
        ->, line width=0.7pt, draw=black!70
    },
    ctrlflow/.style={
        ->, line width=0.7pt, draw=black!55, densely dashed
    },
    edgelbl/.style={
        font=\sffamily\scriptsize, fill=white, inner sep=1.5pt,
        text=black!85, align=center
    }
]

% =========================================================================
% WARSTWA 0: STEROWANIE, EKSPERYMENTY I TESTY (DRY)
% =========================================================================
\node[ctrlbox, text width=2.45cm] (init) at (0, 0) {%
    \texttt{init\_paths.m}\\[1pt]
    {\scriptsize konfiguracja ścieżek}
};
\node[ctrlbox, text width=4.0cm, right=5mm of init] (scripts) {%
    \texttt{scripts/run\_random\_demo.m}\\
    \texttt{scripts/figures/*}\\[1pt]
    {\scriptsize generacja rysunków pracy}
};
\node[ctrlbox, text width=4.1cm, right=5mm of scripts] (mc) {%
    \texttt{scripts/experiments/*}\\
    {\scriptsize \texttt{run\_detection\_monte\_carlo.m}}\\
    {\scriptsize \texttt{run\_range\_limit\_sweep.m}}
};
\node[ctrlbox, text width=3.55cm, right=5mm of mc] (tests) {%
    \texttt{tests/verify\_regression.m}\\[1pt]
    {\scriptsize \textit{Golden Master} i~testy E2E}
};

\draw[ctrlflow] (init) -- (scripts);

\path (init.north) ++(0, 5.2mm) coordinate (ctrl_top_pad);
\begin{scope}[on background layer]
    \node[ctrlframe, fit=(init)(scripts)(mc)(tests)(ctrl_top_pad),
          inner ysep=5pt, inner xsep=7pt] (ctrl_frame) {};
    \node[ctrlhdr] at ([xshift=5pt, yshift=-3.5mm]ctrl_frame.north west)
        {Warstwa inicjalizacji, eksperymentów i~testów (DRY -- bez duplikacji logiki DSP)};
\end{scope}

% =========================================================================
% ETAP 1: NADAJNIK (LEWA KOLUMNA)
% =========================================================================
\node[libbox, text width=6.1cm, anchor=north west] (gwp)
    at ([xshift=7pt, yshift=-14mm]ctrl_frame.south west) {%
    \texttt{lib/get\_wifi\_params.m}\\[1pt]
    {\scriptsize parametry PHY 802.11a/ax, auto-skalowanie ramki, maska aktywna, luki DC}
};
\node[tbbox, text width=6.1cm, below=5.5mm of gwp] (wtx) {%
    \textbf{[WLAN Toolbox]}\\[1pt]
    {\scriptsize \texttt{wlanNonHTConfig} / \texttt{wlanHESUConfig}, \texttt{wlanHEOFDMInfo},}\\
    {\scriptsize \texttt{wlanFieldIndices}, \texttt{wlanSampleRate}, \texttt{transmitTime}, \texttt{wlanWaveformGenerator}}
};

\draw[internalflow] (gwp) -- (wtx);

\path (gwp.north) ++(0, 5.2mm) coordinate (s1_hdr_pad);
\begin{scope}[on background layer]
    \node[stageframe, fit=(gwp)(wtx)(s1_hdr_pad), inner xsep=6pt, inner ysep=5pt] (s1_frame) {};
    \node[stagehdr] at ([xshift=5pt, yshift=-3.6mm]s1_frame.north west)
        {1. Nadajnik: \texttt{src/step1\_transmitter.m}};
\end{scope}

% =========================================================================
% ETAP 2: KANAL PROPAGACYJNY (PRAWA KOLUMNA)
% =========================================================================
\node[libbox, text width=3.25cm, anchor=north east] (ct)
    at ([xshift=-7pt, yshift=-14mm]ctrl_frame.south east) {%
    \texttt{lib/configure\_targets.m}\\[1pt]
    {\scriptsize opóźnienia, Doppler, RCS}\\
    {\scriptsize (\texttt{get\_tgax\_propagation\_params})}
};
\node[libbox, text width=2.85cm, left=3.5mm of ct] (gcf) {%
    \texttt{lib/get\_carrier\_frequency.m}\\[1pt]
    {\scriptsize mapowanie pasma na $f_c$}
};
\node[tbbox, text width=6.75cm, below=5mm of ct.south east, anchor=north east] (wch) {%
    \textbf{[WLAN Toolbox]} \texttt{wlanTGaxChannel} {\scriptsize (Model-A/B/C/E: DPI + clutter)}\\[1pt]
    {\scriptsize + autorska synteza celów ruchomych (opóźnienie ułamkowe FFT + Doppler)}
};
\node[tbbox, text width=6.75cm, below=4.5mm of wch] (ctb) {%
    \textbf{[Communications Toolbox]} \texttt{awgn()}\\[1pt]
    {\scriptsize szum termiczny zakotwiczony względem mocy nominalnej}
};

\draw[internalflow] (gcf.south) -- (gcf.south |- wch.north);
\draw[internalflow] (ct.south) -- (ct.south |- wch.north);
\draw[internalflow] (wch) -- (ctb);

\path (ct.north) ++(0, 5.2mm) coordinate (s2_hdr_pad);
\begin{scope}[on background layer]
    \node[stageframe, fit=(gcf)(ct)(wch)(ctb)(s2_hdr_pad), inner xsep=6pt, inner ysep=5pt] (s2_frame) {};
    \node[stagehdr] at ([xshift=5pt, yshift=-3.6mm]s2_frame.north west)
        {2. Kanał: \texttt{src/step2\_channel.m}};
\end{scope}

% =========================================================================
% ETAP 3: ODBIORNIK - PRZETWARZANIE SYGNALU (PELNA SZEROKOSC)
% =========================================================================
\node[tbbox, text width=4.8cm, anchor=north west] (wsync)
    at ([xshift=7pt, yshift=-14.5mm]s2_frame.south -| ctrl_frame.west) {%
    \textbf{[WLAN Toolbox]} \texttt{wlanPacketDetect}\\
    + \texttt{wlanSymbolTimingEstimate}\\[1pt]
    {\scriptsize + autorski \textit{leading-edge} (L-LTF) / wyzwalacz sprzętowy}
};
\node[libbox, text width=3.4cm, right=5.5mm of wsync] (demod) {%
    \texttt{lib/demodulate.m}\\[1pt]
    {\scriptsize usunięcie CP + kolumnowe FFT}\\
    {\scriptsize (tor Rx i~tor odniesienia Tx)}
};
\node[libbox, text width=5.6cm, right=5.5mm of demod] (est) {%
    \texttt{lib/estimate\_channel.m}\\[1pt]
    {\scriptsize estymacja ZF, filtr MTI ($-\mathrm{mean}$), $\zeta_{\mathrm{ZF}}$}\\
    {\scriptsize wywołuje: \texttt{lib/interpolate\_dc\_nulls.m}}
};

% Drugi rzad Etapu 3 (przeplyw w lewo: est -> rdm -> ax)
\node[libbox, text width=7.5cm, below=6mm of est.south east, anchor=north east] (rdm) {%
    \texttt{lib/compute\_range\_doppler\_map.m}\\[1pt]
    {\scriptsize \textcolor{orange!75!black}{\textbf{[Signal Processing Toolbox]}} \texttt{blackmanharris} (okno 2D)}\\
    {\scriptsize + 2D IFFT/FFT z~dopełnianiem zerami ($4\times$ ZP)}
};
\node[libbox, text width=6.4cm, anchor=west] (ax)
    at (wsync.west |- rdm.center) {%
    \texttt{lib/calculate\_physical\_axes.m}\\[1pt]
    {\scriptsize kalibracja osi fizycznych $(R, v)$ i~złożenie \texttt{radar\_data}}
};

\draw[internalflow] (wsync) -- (demod);
\draw[internalflow] (demod) -- (est);
\draw[internalflow] (est.south) -- (est.south |- rdm.north);
\draw[internalflow] (rdm) -- (ax);

\path (wsync.north) ++(0, 5.2mm) coordinate (s3_hdr_pad);
\begin{scope}[on background layer]
    \node[stageframe, fit=(wsync)(demod)(est)(rdm)(ax)(s3_hdr_pad), inner xsep=6pt, inner ysep=5pt] (s3_frame) {};
    \node[stagehdr] at ([xshift=5pt, yshift=-3.6mm]s3_frame.north west)
        {3. Odbiornik -- przetwarzanie sygnału: \texttt{src/step3\_receiver.m}};
\end{scope}

% =========================================================================
% ETAP 4: ODBIORNIK - DETEKCJA I EKSTRAKCJA (PELNA SZEROKOSC)
% =========================================================================
\node[libbox, text width=3.7cm, anchor=north west] (reg)
    at ([xshift=7pt, yshift=-17mm]s3_frame.south west) {%
    \texttt{lib/get\_detection\_algorithm.m}\\[1pt]
    {\scriptsize centralny rejestr detektorów}\\
    {\scriptsize (\texttt{'binary'} / \texttt{'coherent'})}
};
\node[libbox, text width=6.5cm, anchor=north west] (bstc)
    at ([xshift=10mm, yshift=-11mm]s3_frame.south -| reg.east) {%
    \texttt{lib/binary\_successive\_target\_cancellation.m}\\[1pt]
    {\scriptsize próg CFAR + maskowanie elipsą \textit{null-to-null}}
};
\node[libbox, text width=6.5cm, below=4.5mm of bstc] (cstc) {%
    \texttt{lib/coherent\_successive\_target\_cancellation.m}\\[1pt]
    {\scriptsize próg CFAR + Krok 4 (rezydua) + replika 2D-FFT}
};
\node[libbox, text width=3.1cm, anchor=east] (pnfr)
    at ([xshift=-7pt]s3_frame.east |- reg.center) {%
    \texttt{lib/compute\_pnfr.m}\\[1pt]
    {\scriptsize wyznaczenie kontrastu}\\
    {\scriptsize PNFR wykrytych celów}
};

\draw[internalflow] (reg.east) -- ++(2.5mm, 0) |- (bstc.west)
    node[edgelbl, pos=0.72, above=1pt] {\texttt{'binary'}};
\draw[internalflow] (reg.east) -- ++(2.5mm, 0) |- (cstc.west)
    node[edgelbl, pos=0.72, below=1pt] {\texttt{'coherent'}};
\draw[internalflow] (bstc.east) -- ++(2.5mm, 0) |- (pnfr.west);
\draw[internalflow] (cstc.east) -- ++(2.5mm, 0) |- (pnfr.west);

\path (bstc.north) ++(0, 5.2mm) coordinate (s4_hdr_pad);
\begin{scope}[on background layer]
    \node[stageframe, fit=(reg)(bstc)(cstc)(pnfr)(s4_hdr_pad), inner xsep=6pt, inner ysep=5pt] (s4_frame) {};
    \node[stagehdr] at ([xshift=5pt, yshift=-3.6mm]s4_frame.north west)
        {4. Odbiornik -- detekcja i~ekstrakcja celów: \texttt{src/step4\_interpreter.m}};
\end{scope}

% =========================================================================
% POLACZENIA MIEDZYETAPOWE (PRZEPLYW DANYCH)
% =========================================================================
% Warstwa sterujaca -> Etap 1
\draw[dataflow] ([xshift=1.8cm]ctrl_frame.south west) -- ([xshift=1.8cm]s1_frame.north west)
    node[edgelbl, midway, right=2pt] {parametry scenariusza (\texttt{standard}, \texttt{cbw}, \texttt{mcs}, \texttt{targets})};

% Etap 1 -> Etap 2
\draw[dataflow] (s1_frame.east |- gwp.south) -- (s2_frame.west |- gwp.south)
    node[edgelbl, midway, above=2pt] {\texttt{waveform}, \texttt{fs},\\ \texttt{params}, \texttt{cfg}};

% Etap 1 -> Etap 3 (tor odniesienia Tx w ukladzie wspollokowanym)
\draw[dataflow] (s1_frame.south) -- (s1_frame.south |- s3_frame.north)
    node[edgelbl, midway] {\texttt{waveform}, \texttt{fs}, \texttt{params} (sygnał odniesienia Tx)};

% Etap 2 -> Etap 3 (tor odbiorczy Rx)
\draw[dataflow] (s2_frame.south) -- (s2_frame.south |- s3_frame.north)
    node[edgelbl, midway] {\texttt{signal}, \texttt{fc}, (\texttt{intrinsic\_delay})};

% Etap 3 -> Etap 4
\draw[dataflow] (s3_frame.south) -- (s4_frame.north)
    node[edgelbl, midway] {\texttt{radar\_data} (\texttt{CPer\_F}, \texttt{Per\_F}, \texttt{W}, \texttt{H\_*}, osie $R, v$)};

% Etap 4 -> Warstwa sterujaca / wizualizacyjna (petla powrotna po prawej krawedzi)
\draw[dataflow] (s4_frame.east) -- ++(5.5mm, 0) |- (ctrl_frame.east)
    node[edgelbl, pos=0.26, rotate=90, anchor=south]
        {\texttt{targets}, \texttt{CPer\_F}, \texttt{Per\_F}, \texttt{eta}, \texttt{pnfr}, \texttt{diag\_info}};

% =========================================================================
% LEGENDA
% =========================================================================
\node[anchor=north, yshift=-4.5mm] (legend) at (s4_frame.south) {%
    \begin{tikzpicture}[baseline=-0.5ex, font=\sffamily\scriptsize]
        \node[ctrlframe, minimum width=6.5mm, minimum height=3.5mm, inner sep=0pt] (l1) at (0,0) {};
        \node[right=1.5mm of l1] (t1) {Warstwa sterująca i~testowa (DRY)};
        \node[stageframe, minimum width=6.5mm, minimum height=3.5mm, inner sep=0pt, right=5mm of t1] (l2) {};
        \node[right=1.5mm of l2] (t2) {Fasada potoku (\texttt{src/})};
        \node[libbox, minimum width=6.5mm, minimum height=3.5mm, inner sep=0pt, right=5mm of t2] (l3) {};
        \node[right=1.5mm of l3] (t3) {Kod autorski DSP (\texttt{lib/})};
        \node[tbbox, minimum width=6.5mm, minimum height=3.5mm, inner sep=0pt, right=5mm of t3] (l4) {};
        \node[right=1.5mm of l4] (t4) {Pakiet narzędziowy MATLAB (Toolbox)};
    \end{tikzpicture}
};

\end{tikzpicture}%
}
\caption{Schemat blokowy architektury oprogramowania symulatora radaru pasywnego IEEE~802.11 oraz przepływu danych między modułami autorskimi i~pakietami narzędziowymi MATLAB.}
\label{fig:software_architecture}
\end{figure}
